% from ega0-19-fr.tex
\oldpage[0\textsubscript{IV-1}]{104}

outre, ces corps sont des extensions séparables de $k''$ puisque pour toute
extension $k_1$ de $k''$,
$K\otimes_k k_1=(K\otimes_k k'')\otimes_{k''}k_1$ est sans radical
(\emph{loc. cit.}, th.~1). Si l'on applique ces remarques au corps $K$ de
multiplicité radicielle finie sur $k$, on voit que pour \emph{toute} extension
finie $k''$ de $k_0$, $(K\otimes_k k'')_{\mathrm{red}}$ est séparable sur
$k''$.

Passons à la démonstration de (19.6.6). Il est trivial que $b)$ implique
$b')$. Montrons que $b')$ entraîne $a)$. Si l'on pose
$B'=K'[[T_1,\ldots,T_n]]$, l'hypothèse que $K'$ est séparable sur $k'$
entraîne, par les remarques du début, que pour toute extension finie $k''$ de
$k'$, les composants de l'anneau semi-local complet $B'\otimes_{k'}k''$
(égal à l'anneau de séries formelles
$(K'\otimes_{k'}k'')[[T_1,\ldots,T_n]]$) sont les anneaux
$K_i''[[T_1,\ldots,T_n]]$, où les $K_i''$ sont les corps composants de
$K'\otimes_{k'}k''$. Comme les composants locaux de $B'\otimes_{k'}k''$ sont
aussi des composants locaux de
\[
  \widehat A\otimes_k k''
  =(\widehat A\otimes_k k')\otimes_{k'}k'',
\]
on voit que l'hypothèse $b')$ est encore vérifiée lorsqu'on remplace $k'$ par
une quelconque de ses extensions finies. En particulier, on peut supposer que
$k'$ est \emph{quasi-galoisienne}, donc extension galoisienne d'une extension
radicielle $k'_0$ de $k$~; $A'=\widehat A\otimes_k k'$ peut alors s'écrire
$A'_0\otimes_{k'_0}k'$, où $A'_0=\widehat A\otimes_k k'_0$ est un anneau
local complet. On sait alors (Bourbaki, \emph{Alg.}, chap.~VIII,
\textsection8, prop.~4) que $A'$ est composé direct d'anneaux locaux
\emph{isomorphes} à $A'_0$ en tant que $k'_0$-algèbres. Or, un de ces anneaux
est par hypothèse $k'$-isomorphe à un anneau de séries formelles
$K'[[T_1,\ldots,T_n]]$, où $K'$ est une extension séparable de $k'$, donc
\emph{aussi} une extension séparable de $k'_0$~; compte tenu de (19.6.4),
$A'_0$ est donc une $k'_0$-algèbre formellement lisse. Mais comme $k'_0$ est
un $k$-module fidèlement plat, projectif et de type fini, on conclut de
(19.4.7) que $\widehat A$ est une $k$-algèbre formellement lisse.

On a déjà vu (19.6.5) que $a)$ implique $c)$~; il reste donc à vérifier que
$c)$ implique $b)$. Or, si $c)$ est vérifiée, alors, pour toute extension
radicielle finie $k'$ de $k$ contenant $k_0$, $A\otimes_k k'$ est un anneau
local régulier, dont le corps résiduel $K'$ est une extension séparable de
$k'$~; il résulte donc de (19.6.4) que son complété
$\widehat A\otimes_k k'$ est $k'$-isomorphe à un anneau de séries formelles
$K'[[T_1,\ldots,T_n]]$.
\end{proof}

\begin{remarks}[19.6.7]
\label{0.19.6.7-fr}
\begin{enumerate}
  \item[\rm(i)] On notera que l'hypothèse que $K$ est de multiplicité
  radicielle finie sur $k$ n'a été utilisée que dans la démonstration de
  l'implication $b')\Rightarrow a)$.
  \item[\rm(ii)] Nous prouverons plus tard (22.5.8) que $a)$ et $c)$ sont
  équivalentes, \emph{sans hypothèse} sur l'extension $K$ de $k$.
\end{enumerate}
\end{remarks}

\subsection{Cas des homomorphismes locaux; théorèmes d'existence et
d'unicité}
\label{subsection:0.19.7-fr}

Dans ce numéro, lorsqu'un anneau semi-local est considéré comme un anneau
topologique, il est toujours sous-entendu qu'il s'agit de sa topologie
$\mathfrak r$-préadique, où $\mathfrak r$ est son radical. Tout homomorphisme
local d'anneaux locaux est donc automatiquement \emph{continu}.

\begin{theorem}[19.7.1]
\label{0.19.7.1-fr}
Soient $A$, $B$ deux anneaux locaux noethériens, $\mathfrak m$, $\mathfrak n$
leurs idéaux maximaux respectifs, $k=A/\mathfrak m$ le corps résiduel de $A$~;
on suppose $A$ et $B$ munis respectivement des topologies
$\mathfrak m$-préadique et $\mathfrak n$-préadique. Soit
$\varphi:A\to B$ un homomorphisme local, et posons $B_0=B\otimes_A k$.
Les propriétés suivantes sont équivalentes~:
\begin{enumerate}
  \item[a)] $B$ est une $A$-algèbre \emph{formellement lisse}.

